\documentclass[UTF8, a4paper, fontset=none]{ctexart}

\usepackage[a4paper, margin=1cm, includeheadfoot]{geometry}
\usepackage[most]{tcolorbox}
\usepackage{multicol}

\setCJKmainfont[
	UprightFont	=	NotoSerifTC-Regular.ttf,
	BoldFont	=	NotoSerifTC-Bold.ttf,
	ItalicFont	=	NotoSerifTC-Light.ttf
]{NotoSerifTC-Regular.ttf}
\setCJKsansfont{Noto Sans TC}

\begin{document}
\zihao{-5}

\part{解題導向技巧整理}

\begin{multicols}{3}

\section{高階（常）微分方程（Chapter 4）} 	

\subsection{降階法求齊次解 （§4-2）}

利用降階法，可通過\textbf{已知的一解}，推得另外一解。

\begin{tcolorbox}[size=small, title={\textbf{降階法}（已知一齊次解，待求另一解。）}]
若有一在基本形式 (standard form) 的線性齊次微分方程
$$
y''+P(x)y'+Q(x)y=0
$$
與其一解 $y_1$，則另一解為
$$
y_2 = y_1(x) \int {\frac{e^{-\int{P(x)dx}}}{y_1^2(x)} dx}
$$
\end{tcolorbox}

值得注意的是，若是有超越函數存在於積分之中，則可以
$$
y_2 = y_1(x) \int_{x_0}^{x} {\frac{e^{-\int{P(t)dt}}}{y_1^2(t)} dt}
$$
表達第二解，其中我們假定被積函數在閉區間 $[x_0, x]$ 連續。

\subsection{常係數齊次線性微分方程 (§4-2)}
對於常係數齊次線性微分方程而言，可通過利用\textbf{特徵方程式}將問題化為\textbf{簡單代數問題}，再用解得的代數 $m$ 來推斷解的情形。
\begin{tcolorbox}[size=small, title={\textbf{特徵方程式} (auxiliary equation) }]
對 $n$ 階常係數線性齊次微分方程
$$
a_ny^{(n)}+a_{n-1}y^{(n-1)}+\cdots+a_0y=0
$$
而言，我們可以通過其\textbf{特徵方程式}
$$
a_nm^n+a_{n-1}m^{n-1}+\cdots+a_1m+a_0=0
$$
的解 $m$ 求得對應的微分方程解。解的情形如下所示：
\begin{itemize}
\item \textbf{實根}：該 $m_i$ 對應的解項為 $c_ie^{m_ix}$。
\item \textbf{重根}：該 $m_i$ 若為 $n$ 重根，則對應的解項為 $c_1e^{m_ix}+c_2xe^{m_ix}+\cdots+c_nx^n{m_ix}$。
\item \textbf{虛根}：該 $m_i$ 為兩共軛虛根 $\alpha\pm i\beta$時，其解項依據尤拉公式可得為 $e^{\alpha x} (c_1 \cos{\beta x}+c_2 \sin{\beta x})$。
\end{itemize}
\end{tcolorbox}
此類方法中，我們解特徵方程需要使用到大量的\textbf{綜合除法}、 \textbf{十字交乘}及\textbf{公式解}。請熟練。
\begin{tcolorbox}[size=small, title={一元二次方程的\textbf{公式解}}]
對於一元二次方程式
$$
ax^2+bx+c=0
$$
，其通解為
$$
x=\frac{-b\pm\sqrt{b^2-4ac}}{2a}
$$
\end{tcolorbox}



\end{multicols}

\end{document}
